\documentclass[11pt]{article}

\usepackage[margin=1in]{geometry}
\usepackage{amsmath,amssymb,amsthm}
\usepackage{graphicx}
\usepackage{hyperref}

\newtheorem{theorem}{Theorem}
\newtheorem{lemma}{Lemma}
\newtheorem{fact}{Fact}
\newtheorem{corollary}{Corollary}

\newcommand{\N}{\mathbb N}
\newcommand{\R}{\mathbb R}
\newcommand{\eps}{\varepsilon}
\newcommand{\supp}{\operatorname{supp}}

\title{A slow-weight obstruction to analytic renorming of uncountable Lorentz preduals}
\author{Run \texttt{fa\_banach\_001}, lane 7}
\date{2026-06-20}

\begin{document}
\maketitle

\begin{abstract}
Bible asks in Remark 3.10 of \cite{bible} whether the uncountable Lorentz
preduals \(d_*(w,1,A)\), already shown there to admit \(C^\infty\)-smooth
renormings, admit analytic renormings.  We prove a negative partial result.
Under the standard monotone Lorentz-weight convention, if
\[
  w\notin \ell_q \quad\hbox{for every }1\le q<\infty,
\]
and \(A\) is uncountable, then \(d_*(w,1,A)\) admits no equivalent analytic
norm.  The arbitrary-weight case, in particular weights belonging to some
finite \(\ell_q\), remains open here.
\end{abstract}

\section{Source question}

The source is V. Bible, \emph{Using boundaries to find smooth norms},
arXiv:1311.1408.  The relevant passage is Example 3.9 and Remark 3.10 on
page 15:

\begin{center}
\includegraphics[width=0.94\linewidth]{figures/open_problem_crop.png}
\end{center}

Remark 3.10 says that, for uncountable \(A\), \(X=d_*(w,1,A)\) is a new
example of a space with a \(C^\infty\)-smooth renorming and that it is not yet
known whether \(X\) has an analytic renorming.

\section{Definitions and scope}

Let \(w=(w_j)_{j\ge0}\) be a standard Lorentz weight: \(w_0=1\), \(w_j>0\),
\(w_j\downarrow0\), and \(\sum_j w_j=\infty\).  Put
\[
  W_k=\sum_{j=0}^{k-1} w_j \qquad (k\ge1).
\]
For a set \(A\), \(d_*(w,1,A)\) consists of all scalar functions \(x\) on \(A\)
such that
\[
  \overline x_k
  =
  \sup\left\{
  \frac{\sum_{i=0}^{k-1}|x(a_i)|}{W_k}:
  a_0,\ldots,a_{k-1}\in A\hbox{ distinct}
  \right\}
\]
defines a sequence \((\overline x_k)\in c_0\).  The norm is
\(\|x\|_*=\sup_k \overline x_k\).  Equivalently, if
\((x_i^*)\) is the decreasing rearrangement of the nonzero values of \(|x|\),
then \(\|x\|_*=\sup_k(\sum_{i<k}x_i^*)/W_k\).

The additional hypothesis in this packet is
\[
  \tag{SW} \sum_{j=0}^\infty w_j^q=\infty
  \quad\hbox{for every finite }q\ge1.
\]
This is a genuine restriction.  It covers very slowly decreasing weights, for
instance \(w_j=1/\log(j+2)\), but not weights such as \(w_j=1/(j+1)\).

\section{Proof intuition}

The proof isolates an obstruction to analytic norms that is familiar from
\(c_0(\omega_1)\), but adapted to the Lorentz-predual norm.  Under (SW),
continuous homogeneous polynomials on \(d_*(w,1,A)\) cannot see uncountably
many coordinates.  If an \(n\)-linear form had uncountably many nonzero
coordinate coefficients, a \(\Delta\)-system argument would produce disjoint
coordinate patterns on which the form is uniformly nonzero.  Feeding the form
the vectors \(\sum_j \pm w_j e_{a_j}\), whose \(d_*(w,1,A)\)-norms are bounded
by one, and averaging signs makes the value dominate
\(\sum_j w_j^r\) for some \(1\le r\le n\).  Hypothesis (SW) makes this
impossible.

An analytic function is locally a countable sum of homogeneous polynomials.
Thus, near the origin, it depends on only countably many coordinates.  An
analytic norm would yield an analytic separating function near the origin; a
fresh coordinate outside that countable set gives the contradiction.

\section{The partial result}

\begin{lemma}[Coordinate support of multilinear forms]\label{lem:multi}
Assume (SW), and let \(A\) be uncountable.  Let
\(T\colon X^n\to\R\) be a continuous \(n\)-linear form on
\(X=d_*(w,1,A)\).  Then there is a countable set \(C\subseteq A\) such that
\[
  T(e_{a_1},\ldots,e_{a_n})=0
\]
whenever \(\{a_1,\ldots,a_n\}\not\subseteq C\).
\end{lemma}

\begin{proof}
Suppose not.  By transfinite induction we can choose uncountably many ordered
\(n\)-tuples \(\mathbf a^\xi=(a_1^\xi,\ldots,a_n^\xi)\), \(\xi<\omega_1\),
such that the finite set \(F_\xi=\{a_1^\xi,\ldots,a_n^\xi\}\) is not contained
in the union of its predecessors and
\[
  T(e_{a_1^\xi},\ldots,e_{a_n^\xi})\ne0.
\]
Passing to an uncountable subfamily, assume the absolute values are at least
\(\eps>0\), and that the equality pattern inside the ordered tuples is fixed.
The \(\Delta\)-system lemma gives a further uncountable subfamily for which
the finite sets \(F_\xi\) have common root \(R\), and the petals
\(F_\xi\setminus R\) are pairwise disjoint.  The petals are nonempty by the
choice of the tuples.

Let \(r\) be the number of argument positions occupied by petal coordinates,
counting multiplicity.  Thus \(1\le r\le n\).  Choose any finite subfamily
\(\xi_0,\ldots,\xi_m\).  For root positions use the fixed unit vector
corresponding to that root coordinate.  For a petal position \(i\), consider
vectors of the form
\[
  x_i=\sum_{j=0}^m \eta_{i,j} w_j e_{a_i^{\xi_j}},
  \qquad |\eta_{i,j}|=1.
\]
Because the petal coordinates are distinct and \(w\) is decreasing, every such
\(x_i\) has \(\|x_i\|_*\le1\).

We choose the signs \(\eta_{i,j}\) by the standard random-sign averaging
trick: for each \(j\), choose all but one petal-position signs independently
and choose the last so that the product over the petal positions has the sign
of \(T(e_{a_1^{\xi_j}},\ldots,e_{a_n^{\xi_j}})\).  Averaging over these signs,
all mixed terms with petal indices not all equal have expectation zero, while
the diagonal petal terms contribute positively.  Hence some choice of signs
satisfies
\[
  |T(x_1,\ldots,x_n)|
  \ge
  \eps\sum_{j=0}^m w_j^r .
\]
The left hand side is bounded by \(\|T\|\), since all root and petal input
vectors have norm at most one.  But (SW) implies
\(\sum_j w_j^r=\infty\), a contradiction as \(m\to\infty\).
\end{proof}

\begin{corollary}[Coordinate support of polynomials]\label{cor:poly}
Assume (SW), and let \(A\) be uncountable.  Every continuous homogeneous
polynomial on \(d_*(w,1,A)\) is countably supported: for each such polynomial
\(P\) there is a countable \(C\subseteq A\) such that
\[
  P(x)=P(x\mathbf 1_C)\qquad (x\in d_*(w,1,A)).
\]
\end{corollary}

\begin{proof}
Apply Lemma \ref{lem:multi} to the symmetric multilinear form associated to
\(P\), and use the polarization formula.
\end{proof}

\begin{fact}[analytic separation]\label{fact:sep}
If a normed space admits an equivalent analytic norm, then there are
\(\rho>0\) and an analytic function \(f\), represented by a power series at
the origin on a ball containing \(\{x:\|x\|=\rho\}\), such that
\[
  f(0)=0
  \quad\hbox{and}\quad
  \inf_{\|x\|=\rho} f(x)>1.
\]
\end{fact}

\noindent This is the standard separating-function consequence of an analytic
norm, used for example in the proof that no dense subspace of
\(c_0(\omega_1)\) admits an analytic norm in \cite{dantas-hajek-russo}.

\begin{theorem}\label{thm:main}
Let \(A\) be uncountable and let \(w\) satisfy (SW).  Then
\[
  X=d_*(w,1,A)
\]
does not admit an equivalent analytic norm.
\end{theorem}

\begin{proof}
Assume that \(X\) admits an equivalent analytic norm.  Since equivalent norms
give the same continuous polynomials, Corollary \ref{cor:poly} applies to the
homogeneous polynomials in the local power-series expansion of any analytic
function for that norm.

Let \(f\) and \(\rho\) be as in Fact \ref{fact:sep}, and write the expansion
of \(f\) at zero as
\[
  f(x)=\sum_{n=0}^\infty P_n(x)
\]
on a ball containing the \(\rho\)-sphere.  By Corollary \ref{cor:poly}, each
\(P_n\) is supported on a countable set \(C_n\subseteq A\).  Let
\(C=\bigcup_n C_n\), still countable.  Choose \(a\in A\setminus C\) and take
\[
  x=\rho\, e_a/\|e_a\|
\]
with respect to the analytic norm.  Then \(x\) lies on the \(\rho\)-sphere,
but \(x\mathbf 1_C=0\).  Therefore
\[
  f(x)=\sum_n P_n(x)=\sum_n P_n(0)=f(0)=0,
\]
contradicting \(\inf_{\|y\|=\rho}f(y)>1\).
\end{proof}

\section{What remains open}

The theorem does not settle Bible's Remark 3.10 for all Lorentz weights.  The
remaining case is when \(w\in\ell_q\) for some finite \(q\).  In that regime
the obstruction above is no longer available: diagonal high-degree polynomials
such as \(\sum_{a\in A}x(a)^q\) can become continuous, so analytic functions
need not be countably supported.

\section{Verification and novelty}

The proof is symbolic and has no computational dependency.  The source crop
was generated from the arXiv PDF by \texttt{code/crop\_open\_problem.py}.

The run indexes were searched for \texttt{1311.1408}, \texttt{Using
boundaries to find smooth norms}, \texttt{d\_*(w,1,A)}, \texttt{Lorentz
predual}, and \texttt{analytic renorming}; no prior packet for this target
was found.  Web searches on 2026-06-20 for the exact Lorentz-predual analytic
renorming question did not locate a later full answer.  A nearby later paper
\cite{dantas-hajek-russo} supplies the model analytic-support obstruction for
\(c_0(\omega_1)\), but it does not directly answer Bible's Banach-space
question because the \(d_*(w,1,A)\)-norm is stronger than the inherited
\(c_0(A)\)-norm.

\section{Human review recommendation}

Review should focus on the sign-averaging step in Lemma \ref{lem:multi} and
on the monotone-weight convention.  If the source permits non-monotone weights
literally, the packet should be read under the standard Lorentz replacement by
the nonincreasing rearrangement of \(w\), or the theorem statement should keep
the explicit monotonicity assumption.

\begin{thebibliography}{9}

\bibitem{bible}
V. Bible,
\emph{Using boundaries to find smooth norms},
arXiv:1311.1408, 2013.

\bibitem{dantas-hajek-russo}
S. Dantas, P. H\'ajek, and T. Russo,
\emph{Smooth norms in dense subspaces of Banach spaces},
arXiv:1911.09611, 2019.

\end{thebibliography}

\end{document}
